\documentclass{book}
\usepackage[backend=biber,natbib=true,style=authoryear]{biblatex}
\addbibresource{/home/nguyen/1_NQBH/math/bib.bib}
\usepackage[utf8]{inputenc}
\usepackage{float}
\usepackage{graphicx}
\usepackage[colorlinks=true,linkcolor=blue,urlcolor=red,citecolor=magenta]{hyperref}
\usepackage{amsmath,amssymb,amsthm}
\allowdisplaybreaks
\numberwithin{equation}{section}
\newtheorem{assumption}{Assumption}[section]
\newtheorem{lemma}{Lemma}[section]
\newtheorem{corollary}{Corollary}[section]
\newtheorem{definition}{Definition}[section]
\newtheorem{proposition}{Proposition}[section]
\newtheorem{theorem}{Theorem}[section]
\newtheorem{notation}{Notation}[section]
\newtheorem{remark}{Remark}[section]
\newtheorem{example}{Example}[section]
\newtheorem{ques}{Question}[section]
\newtheorem{problem}{Problem}[section]
\newtheorem{conjecture}{Conjecture}[section]
\usepackage[left=0.5in,right=0.5in,top=1.5cm,bottom=1.5cm]{geometry}
\usepackage{fancyhdr}
\pagestyle{fancy}
\fancyhf{}
\addtolength{\headheight}{0pt}% obsolete
\lhead{\scriptsize \chaptername~\thechapter}
\rhead{\itshape \scriptsize \nouppercase{\leftmark}} %\nouppercase !
\renewcommand{\chaptermark}[1]{\markboth{#1}{}}
\cfoot{\thepage}

\title{William Strunk Jr., E.B. White (2019). The Elements of Style., 4th Edition}
\author{Collector: Nguyen Quan Ba Hong}
\date{\today}

\begin{document}
\maketitle
\setcounter{secnumdepth}{6}
\setcounter{tocdepth}{6}
\tableofcontents

%------------------------------------------------------------------------------%

\section*{Amazon/reviews}
\begin{quotation}
    `The Elements of Style' (1918), by William Strunk, Jr., is an American English writing style guide.
    
    It is the best-known, most influential prescriptive treatment of English grammar and usage, and often is required reading and usage in U.S. high school and university composition classes.
    
    This edition of `The Elements of Style' details 8 elementary rules of usage, 10 elementary principles of composition, ``a few matters of form'', and a list of commonly misused words and expressions.
\end{quotation}

\begin{quotation}
    ``\textit{$\ldots$ a marvellous and timeless little book$\ldots$ Here, succinctly, elegantly and without fuss are the essentials of writing clear, correct English}.'' - John Clare, \textit{The Telegraph}
\end{quotation}

%------------------------------------------------------------------------------%

\chapter*{Forword}

The 1st writer Roger Angell watched at work was my stepfather, E. B. White.

Each Tuesday morning, he would close his study door and sit down to write the ``Notes \& Comment'' page for \textit{The New Yorker}.

The task was familiar with him - he was required to file a few hundred words of editorial or personal commentary on some topic in or out of the news that week - but the sounds of his typewriter from his room came in hesitant bursts, with long silences in between.

Hours went by.

Summoned at last for lunch, he was silent and preoccupied, and soon excused himself to get back to the job.

When the copy went off at last, in the afternoon RFD pouch - we were in Maine, a day's mail away from New York - he rarely seemed satisfied.

``It isn't good enough,'' he said sometimes.

``I wish it were better.''

%
Writing is hard, even for authors who do it all the time.

Less frequent practitioners - the job applicant; the business executive with an annual report to get out; the high school senior with a Faulkner assignment; the graduate-school student with her thesis proposal; the writer of a letter of condolence - often get stuck in an awkward passage of find a muddle on their screens, and then blame themselves.

What should be easy and flowing looks tangled or feeable or overblown - not what was meant at all.

What's wrong with me, each one thinks.

Why can't I get this right?

%
It was this recurring question, put to himself, that must have inspired White to revive and add to a textbook by an English professor of his, Will Strunk Jr., that he had 1st read in college, and to get it published.

The result, this quiet book, has been in print for 40 years, and has offered more than 10 million writers a helping hand.

White knew that a compendium of specific tips - about singular and plural verbs, parentheses, that ``that'' - ``which'' scuffle, and many others - could clear up a recalcitrant sentence or subclause when quickly reconsulted, and that the larger principles needed to be kept in plain sight, like a wall sampler.

%
How simple they look, set down here in White's last chapter: ``Write in a way that comes naturally,'' ``Revise and rewrite,'' ``Do not explain too much,'' and the rest; above all, the cleansing, clarion ``Be clear.''

How often Roger Angell has turned to them, in the book or in my mind, while trying to start or unblock or revise some piece of my own writing!

They help - they really do.

They work.

They are the way.

%
E. B. White's prose is celebrated for its ease and clarity - just think of \textit{Charlotte's Web} - but maintaining this standard required endless attention.

When the new issue of \textit{The New Yorker} turned up in Maine, Roger Angell sometimes saw him reading his ``Comment'' piece over to himself, with only a slightly different expression than the one he'd worn on the day it went off.

Well, O.K., he seemed to be saying.

\textit{At least I got the elements right}.

%
This edition has been modestly updated, with word processors and air conditioners making their 1st appearance among White's references, and with a light redistribution of genders to permit a feminine pronoun or female farmer to take their places among the males who once innocently served him.

Sylvia Plath has knocked Keats out of the box, and Roger Angell notices that ``America'' has become ``this country'' in a sample text, to forestall a subsequent and possibly demeaning ``she'' in the same paragraph.

What is not here is anything about E-mail - the rules-free, lower-case flow that cheerfully keeps us in touch these days.

E-mail is conversation, and it may be replacing the sweet and endless talking we once sustained (and tucked away) within the informal letter.

But we are all writers and readers as well as communicators, with the need at times to please and satisfy ourselves (as White put it) with the clear and almost perfect thought.

\begin{flushright}
    \textit{Roger Angell}
\end{flushright}

%------------------------------------------------------------------------------%

\chapter*{Introduction}

At the close of the 1st World War, when E. B. White was a student at Cornell, E. B. White took a course called \textit{English 8}.

My professor was William Strunk Jr.

A textbook required for the course was a slim volume called \textit{The Elements of Style}, whose authors was the professor himself.

The year was 1919.

The book was known on the campus in those days as ``the little book,'' with the stress on the word ``little.''

It had been privately printed by the author.

%
E. B. White passed the course, graduated from the university, and forgot the book but not the professor.

Some 38 years later, the book bobbed up again in my life when Macmillan commissioned me to revise it for the college market and the general trade.

Meantime, Prof. Strunk had died.

%
\textit{The Elements of Style}, when E. B. White reexamined it in 1957, seemed to me to contain rich deposits of gold.

It was Will Strunk's \textit{parvum opus}, his attempt to cut the vast tangle of English rhetoric down to size and write its rules and principles on the head of a pint.

Will himself had hung the tag ``little'' on the book; he referred to it sardonically and with secret pride as ``the \textit{little} book,'' always giving the word ``little'' a special twist, as though he were putting a spin on a ball.

In its original form, it was a 43 page summation of the case for cleanliness, accuracy, and brevity in the use of English.

Today, 52 years later, its vigor is unimpaired, and for sheer pith E. B. White thinks it probably sets a record that is not likely to be broken.

Even after E. B. White got through tampering with it, it was still a tiny thing, a barely tarnished gem.

7 rules of usage, 11 principles of composition, a few matters of form, and a list of words and expressions commonly misused - that was the sum and substance of Prof. Strunk's work.

Somewhat audaciously, and in an attempt to give my publisher his money's worth, E. B. White added a chapter called ``An Approach to Style,'' setting forth my own prejudices, my notions of error, my articles of faith.

This chapter (Chap. V) is addressed particularly to those who feel that English prose composition is not only a necessary skills but a sensible pursuit as well - a way to spend one's days.

E. B. White thinks Prof. Strunk would not object to that.

%
A 2nd edition of the book was published in 1972.

E. B. White has now completed a 3rd revision.

Chap. IV has been refurbished with words and expressions of a recent vintage; 4 rules of usage have been added to Chap. I.

Fresh examples have been added to some of the rules and principles, amplification has reared its head in a few places in the text where E. B. White felt an assault could successfully be made on the bastions of its brevity, and in general the book has received a thorough overhaul - to correct errors, delete bewhiskered entries, and enliven the argument.

%
Prof. Strunk was a positive man.

His book contains rules of grammar phrased as direct orders.

In the main E. B. White has not tried to soften his commands, or modify his pronouncements, or remove the special objects of his scorn.

E. B. White has tried, instead, to preserve the flavor of his discontent while slightly enlarging the scope of the discussion.

\textit{The Element of Style} does not pretend to survey the whole field.

Rather it proposes to give in brief space the principal requirements of plain English style.

It concentrates on fundamentals: the rules of usage and principles of composition most commonly violated.

%
The reader will soon discover that these rules and principles are in the form of sharp commands, Sergeant Strunk snapping orders to his platoon.

``Do not joint independent clauses with a comma.'' (Rule 5.)

``Do not break sentences in 2.'' (Rule 6.)

``Use the active voice.'' (Rule 14.)

``Omit needless words.'' (Rule 17.)

``Avoid a succession of loose sentences.'' (Rule 18.)

``In summaries, keep to 1 tense.'' (Rule 21.)

Each rule or principle is followed by a short hortatory essay, and usually the exhortation is followed by, or interlarded with, examples in parallel columns - the true vs. the false, the right vs. the wrong, the timid vs. the bold, the ragged vs. the trim.

From every line there peers out at me the puckish face of my professor, his short hair pared neatly in the middle and combed down over his forehead, his eyes blinking incessantly behind steel-rimmed spectacles as though he had just emerged into strong light, his lips nibbling each other like nervous horses, his smile shuttling to and fro under a carefully edged mustache.

%
``Omit needless words!'' cries the author on p. 24, and into that imperative Will Strunk really put his heart and soul.

In the days when E. B. White was sitting in his class, he omitted so many needless words, and omitted them so forcibly and with such eagerness and obvious relish, that he often seemed in the position of having shortchanged himself - a man left with nothing more to say yet with time to fill, a radio prophet who had out-distanced the clock.

Will Strunk got out of this predicament by a simple trick: he uttered every sentence 3 times.

When he delivered his oration on brevity to the class, he leaned forward over his desk, grasped his coat lapels in his hands, and, in a husky, conspiratorial voice, said, ``Rule 17. Omit needless words! Omit needless words! Omit needless words!''

%
He was a memorable man, friendly and funny.

Under the remembered string of his kindly lash, E. B. White has been trying to omit needless words since 1919, and although there are still many words that cry for omission and the huge task will never be accomplished, it is exciting to me to reread the masterly Strunkian elaboration of this noble theme.

It goes:

\begin{quotation}
    \it
    Vigorous writing is concise. A sentence should contain no unnecessary words, a paragraph no unnecessary sentences, for the same reason that a drawing should have no unnecessary lines and a machine no unnecessary parts.
    
    This requires not that the writer make all sentences short or avoid all detail and treat subjects only in outline, but that every word tell.
\end{quotation}
There you have a short, valuable essay on the nature and beauty of brevity - 59 words that could change the world.

Having recovered from his adventure in prolixity (59 words were a lot of words in the tight world of William Strunk Jr.), the professor proceeds to give a few quick lessons in pruning.

Students learn to cut the dead-wood from ``this is a subject that,'' reducing it to ``this subject,'' a saving of 3 words.

They learn to trim ``used for fuel purposes'' down to ``used for fuel.''

They learn that they are being chatterboxes when they say ``the question as to whether'' and they should just day ``whether'' - a saving of 4 words out of a possible 5.

%
The professor devotes a special paragraph to the vile expression \textit{the fact that}, a phrase that causes him to quiver with revulsion.

The expression, he says, should be ``revised out of every sentence in which it occurs.''

But a shadow of gloom seems to hang over the page, and you feel that he knows how hopeless his cause is.

E. B. White supposes E. B. White has written \textit{the fact that} a thousand times in the heat of composition, revised it out maybe 500 times in the cool aftermath.

To be batting only .500 this late in the season, to fail half the time to connect with this fat pitch, saddens me, for it seems a betrayal of the man who showed me how to swing at it and made the swinging seem worthwhile.

%
E. B. White treasures \textit{The Elements of Style} for its shape advice, but E. B. White treasures it even more for the audacity and self-confidence of its author.

Will knew where he stood.

He was so sure of where he stood, and made his position so clear and so plausible, that his peculiar stance has continued to invigorate me - and, E. B. White is sure, thousands of other ex-students - during the years that have intervened since our 1st encounter.

He had a number of likes and dislikes that were almost as whimsical as the choice of a necktie, yet he made them seem utterly convincing.

He disliked the word \textit{forceful} and advised us to use \textit{forcible} instead.

He felt that the word \textit{clever} was greatly overused: ``It is best restricted to ingenuity displayed in small matters.''

He despised the expression \textit{student body}, which he termed gruesome, and made a special trip downtown to the \textit{Alumni News} office 1 day to protest the expression and suggest that \textit{studentry} be substituted - a coinage of is own, which he felt was similar to \textit{citizenry}.

E. B. White is told that the \textit{News} editor was so charmed by the visit, if not by the word, that he ordered the student body buried, never to rise again.

\textit{Studentry} has taken its place.

It's not much of an improvement, but it does sound less cadaverous, and it made Will Strunk quite happy.

%
Some years ago, when the heir to the throne of England was a child, E. B. White noticed a headline in the \textit{Times} about Bonnie Prince Charlie: ``CHARLES' TONSILS OUT.''

Immediately Rule 1 leapt to mind.
\begin{enumerate}
    \item Form the posssessive singular of nuns by adding \textit{'s}.
    
    Follow this rule whatever the final consonant.
    
    Thus write,
    \begin{example}
        Charles's friend
        
        Burns's poems
        
        the witch's malice
    \end{example}
\end{enumerate}
Clearly, Will Strunk had foreseen, as far back as 1918, the dangerous tonsillectomy of a prince, in which the surgeon removes the tonsils and the \textit{Times} copy desk removes the final \textit{s}.

He started his book with it.

E. B. White recommends Rule 1 to the \textit{Times}, and E. B. White trusts that Charles's throat, not Charles' throat, is in fine shape today.

%
Style rules of this sort are, of course, somewhat a matter of individual preference, and even the established rules of grammar are open to challenge.

Professor Strunk, although 1 of the most inflexible and choosy of men, was quick to acknowledge the fallacy of inflexibility and the danger of doctrine.
\begin{quotation}
    ``\textit{It is an old observation},'' he wrote, ``\textit{that the best writers sometimes disregard the rules of rhetoric. When they do so, however, the reader will usually find in the sentence some compensating merit, attained at the cost of the violation. Unless he is certain of doing as well, he will probably do best to follow the rules}.''
\end{quotation}
It is encouraging to see how perfectly a book, even a dusty rule book, perpetuates and extends the spirit of a man.

Will Strunk loved the clear, the brief, the bold, and his book is clear, brief, bold.

Boldness is perhaps its chief distinguishing mark.

On p. 26, explaining 1 of his parallels, he says,
\begin{quotation}
    ``\textit{The lefthand version gives the impression that the writer is undecided or timid, apparently unable or afraid to choose 1 form of expression and hold to it}.''
\end{quotation}
And his original Rule 11 was ``Make definite assertions.''

That was Will all over.

He scorned the vague, the tame, the colorless, the irresolute.

He felt it was worse to be irresolute than to be wrong.

E. B. White remembers a day in class when he leaned far forward, in this characteristic pose - the pose of a man about to impart a secret - and croaked,
\begin{quotation}
    ``\textit{If you don't know how to pronounce a word, say it loud! If you don't know how to pronounce a word, say it loud!}''
\end{quotation}
This comical piece of advice struck me as sound at the time, and E. B. White still respects it.

\textit{Why compound ignorance with inaudibility?}

\textit{Why run and hide?}

%
All through \textit{The Elements of Style} one finds evidences of the author's deep sympathy for the reader.

Will felt that the reader was in serious trouble most of the time, floundering in a swamp, and that it was the duty of anyone attempting to write English to drain this swamp quickly and get the reader up on dry ground, or at least to throw a rope.

In revising the text, E. B. White has tried to hold steadily in mind this belief of his, this concern for the bewildered reader.

%
In the English classes of today, ``the little book'' is surrounded be longer, lower textbooks - books with permissive steering and automatic transitions.

Perhaps the book has become something of a curiosity.

To E. B. White, it still seems to maintain its original poise, standing, in a drafty time, erect, resolute, and assured.

E. B. White still finds the Strunkian wisdom a comfort, the Strunkian humor a delight, and the Strunkian attitude toward right-and-wrong a blessing undisguised.

\begin{flushright}
    E. B. White
\end{flushright}

%------------------------------------------------------------------------------%

\chapter{Elementary Rules of Usage}

\section{Form the possessive singular of nouns by adding `s}
Follow this rule whatever the final consonant.

Thus write,
\begin{example}
    Charles's friend
    
    Burns's poems
    
    the witch's malice
\end{example}
Exceptions are the possessive of ancient proper names in \textit{-es} an \textit{-is}, the possessive \textit{Jesus'}, and such forms as \textit{for conscience' sake, for righteousness' sake}.

But such forms as \textit{Achilles' heel, Moses' laws, Isis' temple} are commonly replaced by
\begin{example}
    the laws of Moses
    
    the temple of Isis
\end{example}
The pronominal possessives \textit{hers, its, theirs, yours}, and \textit{ours} have no apostrophe.

Indefinite pronouns, however, use the apostrophe to show possession.
\begin{example}
    one's rights
    
    somebody else's umbrella
\end{example}
A common error is to write \textit{it's} for \textit{its}, or vice versa.

The 1st is a contraction, meaning ``it is.''

The 2nd is a possessive.
\begin{example}
    It's wise dog that scratches its own fleas.
\end{example}

\section{In a series of 3 or more terms with a single conjunction, use a comma after each term except the last}
Thus write,
\begin{example}
    red, white, and blue
    
    gold, silver, or copper
    
    He opened the letter, read it, and made a note of its contents.
\end{example}
This comma is often referred to as the ``serial'' comma.

In the names of business firms the last comma is usually omitted.

Follow the usage of the individual firm.
\begin{example}
    Little, Brown and Company
    
    Donaldson, Lufkin \& Jenrette
\end{example}

\section{Enclose parenthetic expressions between commas}
\begin{example}
    The best way to see a country, unless you are pressed for time, is to travel on foot.
\end{example}
This rule is difficult to apply; it is frequently hard to decide whether a single word, e.g., \textit{however}, or a brief phrase is or is not parenthetic.

If the interruption to the flow of the sentence is but slight, the commas may be safely omitted.

But whether the interruption is slight or considerable, never omit 1 comma and leave the other.

There is no defense for such punctuation as
\begin{example}
    Marjories husband, Colonel Nelson paid us a visit yesterday.
\end{example}
or
\begin{example}
    My brother you will be pleased to hear, is now in perfect health.
\end{example}
Dates usually contain parenthetic words or figures.

Punctuate as follows:
\begin{example}
    February to July, 1992
    
    April 6, 1986
    
    Wednesday, November 14, 1990
\end{example}
Note that it is customary to omit the comma in
\begin{example}
    6 April 1988
\end{example}
The last form is an excellent way to write a date; the figures are separated by a word and are, for that reason, quickly grasped.

%
A name or a title in direct address is parenthetic.
\begin{example}
    If, Sir, you refuse, I cannot predict what will happen.
    
    Well, Susan, this is a fine mess you are in.
\end{example}
The abbreviation \textit{etc., i.e.}, and \textit{e.g.}, the abbreviations for academic degrees, and titles that follow a name are parenthetic and should be punctuated accordingly.
\begin{example}
    Letters, packages, etc., should go here.
    
    Horace Fulsome, Ph.D., presided.
    
    Rachel Simonds, Attorney
    
    The Reverend Harry Lang, S.J.
\end{example}
No comma, however, should separate a noun from a restrictive term of identification.
\begin{example}
    Billy the Kid
    
    The novelist Jane Austen
    
    William the Conqueror
    
    The poet Sappho
\end{example}
Although \textit{Junior}, with its abbreviation \textit{Jr.}, has commonly been regarded as parenthetic, logic suggests that it is, in fact, restrictive and therefore not in need of a comma.
\begin{example}
    James Wright Jr.
\end{example}
Nonrestrictive relative clauses are parenthetic, as are similar clauses introduced by conjunctions indicating time or place.

Commas are therefore needed.

A nonrestrictive clauses is one that does not serve to identify or define the antecedent noun.
\begin{example}
    The audience, which had at 1st been indifferent, became more and more interested.
    
    In 1769, when Napoleon was born, Corsica had but recently been acquired by France.
    
    Nether Stowey, where Coleridge wrote The Rime of the Ancient Mariner, is a few miles from Bridgewater.
\end{example}
In these sentences, the clauses introduced by \textit{which, when}, and \textit{where} are nonrestrictive; they do not limit or define, they merely add something.

In the 1st example, the clause introduced by \textit{which} does not serve to tell which of several possible audiences is meant; the reader presumably knows that already.

The clause adds, parenthetically, a statement supplementing that in the main clause.

Each of the 3 sentences is a combination of 2 statements that might have been made independently.
\begin{example}
    The audience was at 1st indifferent. Later it became more and more interested.
    
    Napoleon was born in 1769. At that time Corsica had but recently been acquired by France.
    
    Coleridge wrote The Rime of the Ancient Mariner at Nether Stowey. Nether Stowey is a few miles from Bridgewater.
\end{example}
Restrictive clauses, by contrast, are not parenthetic and not set off by commas.

Thus.
\begin{example}
    People who live in glass houses shouldn't throw stones.
\end{example}
Here the clause introduced by \textit{who} does serve to tell which people are meant; the sentence, unlike the sentences above, cannot be split into 2 independent statements.

The same principle of comma use applies to participial phrases and to appositives.
\begin{example}
    People sitting in the rear couldn't hear, \emph{(restrictive)}
    
    Uncle Bert, being slightly deaf, moved forward, \emph{(non-restrictive)}
    
    My cousin Bob is a talented harpist, \emph{(restrictive)}
    
    Our oldest daughter, Mary, sings, \emph{nonrestrictive}
\end{example}
When the main clause of a sentence is preceded by a phrase or a subordinate clause, use a comma to set off these elements.
\begin{example}
    Partly by hard fighting, partly by diplomatic skill, they enlarged their dominions to the east and rose to royal rank with the possession of Sicily.
\end{example}

\section{Place a comma before a conjunction introduction an independent clause}
\begin{example}
    The early records of the city have disappeared, and the story of its 1st years can no longer be reconstructed.
    
    The situation is perilous, but there is still 1 chance of escape.
\end{example}
2-part sentences of which the 2nd member is introduced by as (in the sense of ``because''), \textit{for, or, nor}, or \textit{while} (in the sense of ``and at the same time'') likewise require a comma before the conjunction.

%
If a dependent clause, or an introductory phrase requiring to be set off by a comma, precedes the 2nd independent clause, no comma is needed after the conjunction.
\begin{example}
    The situation is perilous, but if we are prepared to act promptly, there is still 1 chance of escape.
\end{example}
When the subject is the same for both clauses and is expressed only once, a comma is useful if the connective is \textit{but}.

When the connective is \textit{and}, the comma should be omitted if the relation between the 2 statements is close or immediate.
\begin{example}
    I have heard the arguments, but am still unconvinced.
    
    He has had several year's experience and is thoroughly competent.
\end{example}

\section{Do not join independent clauses with a comma}
If 2 or more clauses grammatically complete and not joined by a conjunction are to form a single compound sentence, the proper mark of punctuation is a semicolon.
\begin{example}
    Mary Shelley's works are entertaining; they are full of engaging ideas.
    
    It is nearly half past 5; we cannot reach town before dark.
\end{example}
It is, of course, equally correct to write each of these as 2 sentences, replacing the semicolons with periods.
\begin{example}
    Mary Shelley's works are entertaining. They are full of engaging ideas.
    
    It is nearly half past 5. We cannot reach town before dark.
\end{example}
If a conjunction is inserted, the proper mark is a comma. (Rule 4.)
\begin{example}
    Mary Shelley's works are entertaining, for they are full of engaging ideas.
    
    It is nearly half past 5, and we cannot reach town before dark.
\end{example}
A comparison of the 3 forms given above will show clearly the advantage of the 1st.

It is, at least in the examples given, better than the 2nd form because it suggests the close relationship between the 2 statements in a way that the 2nd does not attempt, and better than the 3rd because it is briefer and therefore more forcible.

Indeed, this simple method of indicating relationship between statements is 1 of the most useful devices of composition.

The relationship, as above, is commonly 1 of cause and consequence.

%
Note that if the 2nd clause is preceded by an adverb, e.g., \textit{accordingly, besides, then, therefore}, or \textit{thus}, and not by a conjunction, the semicolon is still required.
\begin{example}
    I had never been in the place before; besides, it was dark as a tomb.
\end{example}
An exception to the semicolon rule is worth noting here.

A comma is preferable when the clauses are very short and alike in form, or when the tone of the sentence is easy and conversational.
\begin{example}
    Man proposes, God disposes.
    
    The gates swung apart, the bridge fell, the portcullis was drawn up.
    
    I hardly knew him, he was so changed.
    
    Here today, gone tomorrow.
\end{example}

\section{Do not break sentences in 2}
In other words, do not use periods for commas.
\begin{example}
    I met them on a Cunard liner many years ago. Coming home from Liverpool to New York.
    
    She was an interesting talker. A woman who had traveled all over the world and lived in half of dozen countries.
\end{example}
In both these examples, the 1st period should be replaced by a comma and the following word begun with a small letter.

%
It is permissible to make an emphatic word or expression serve the purpose of a sentence and to punctuate it accordingly:
\begin{example}
    Again and again he called out. No reply.
\end{example}
The writer must, however, be certain that the emphasis is warranted, lest a clipped sentence seem merely a blunder in syntax or in punctuation.

Generally speaking, the place for broken sentences is in dialogue, when a character happens to speak in a clipped or fragmentary way.

%
Rules 3, 4, 5, and 6 cover the most important principles that govern punctuation.

They should be so thoroughly mastered that their application becomes 2nd nature.

\section{Use a colon after an independent clause to introduce a list of particulars, an appositive, an amplification, or an illustrative quotation}
A colon tells the reader that what follows is closely related to the preceding clause.

The colon has more effect than the comma, less power to separate than the semicolon, and more formality than the dash.

It usually follows an independent clause and should not separate a verb from its complement or a preposition from its object.

The examples in the lefthand column, below, are wrong; they should be rewritten as in the righthand column.
\begin{example}
    Your dedicated whittler requires: a knife, a piece of wood, and a back porch.
    
    Understanding is that penetrating quality of knowledge that grows from: theory, practice, conviction, assertion, error, and humiliation.
    
    Your dedicated whittler requires 3 props: a knife, a piece of wood, and a back porch.
    
    Understanding is that penetrating quality of knowledge that grows from theory, practice, conviction, assertion, error, and humiliation.
\end{example}
Join 2 independent clauses with a colon if the 2nd interprets or amplifies the 1st.
\begin{example}
    But even so, there was a directness and dispatch about animal burial: there was no stopover in the undertaker's foul parlor, no wreath or spray.
\end{example}
A colon may introduce a quotation that supports or contributes to the preceding clause.
\begin{example}
    The squalor of the streets reminded her of a line from Oscar Wilde: ``We are all in the gutter, but some of us are looking at the star.''
\end{example}
The colon also has certain functions of form: to follow the salutation of a former letter, to separate hour from minute in a notation of time, and to separate the title of work from its subtitle or a Bible chapter from a verse.
\begin{example}
    Dear Mr. Montague:
    
    departs at 10:48 P.M.
    
    Practical Calligraphy: An Introduction to Italic Script
\end{example}

\section{Use a dash to set off an abrupt break or interruption and to announce a long appositive or summary}
A dash is a mark of separation stronger than a comma, less formal than a colon, and more relaxed than parentheses.
\begin{example}
    His 1st thought on getting out of bed - if he had any thought at all - was to get back in again.
    
    The rear axle began to make a noise - a grinding, chattering, teeth-gritting rasp.
    
    The increasing reluctance of the sun to rise, the extra nip in the breeze, the patter of shed leaves dropping - all the evidences of fall drifting into winter were clearer each day.
\end{example}
Use a dash only when a more common mark of punctuation seems inadequate.
\begin{example}
    Her father's suspicions proved well-founded - it was not Edward she cared for - it was San Francisco.
    
    $\to$ Her father's suspicions proved well-founded. It was not Edward she cared for, it was San Francisco.
    
    Violence - the kind you see on television - is not honestly violent - there lies its harm.
    
    $\to$ Violence, the kind you see on television, is not honestly violent. There lies its harm.
\end{example}

\section{The number of the subject determines the number of the verb}
Words that intervene between subject and verb do not affect the number of the verb.
\begin{example}
    The bittersweet flavor of youth - its trials, its joys, its adventures, its challenges - are not soon forgotten.
    
    $\to$ The bittersweet flavor of youth - its trials, its joys, its adventures, its challenges - is not soon forgotten.
\end{example}
A common blunder is the use of a singular verb form in a relative clause following ``one of$\ldots$'' or a similar expression when the relative is the subject.
\begin{example}
    1 of the the ablest scientists who has attacked this problem.
    
    $\to$ 1 of the ablest scientists who have attached this problem.
    
    1 of those people who is never ready on time $\to$ 1 of those people who are never ready on time
\end{example}
Use a singular verb form after \textit{each, either, everyone, everybody, neither, nobody, someone}.
\begin{example}
    Everybody thinks he has a unique sense of humor.
    
    Although both clocks strike cheerfully, neither keeps good time.
\end{example}
With \textit{none}, use the singular verb when the word means ``no one'' or ``not one.''
\begin{example}
    None of us are perfect.
    
    None of us is perfect.
\end{example}
A plural verb is commonly used when \textit{none} suggests more than 1 thing or person.
\begin{example}
    None are so fallible as those who are sure they're right.
\end{example}
A compound subject formed of 2 or more nouns joined by \textit{and} almost always requires a plural verb.
\begin{example}
    The walrus and the carpenter were walking close at hand.
\end{example}
But certain compounds, often cliches, are so inseparable they are considered a unit and so take a singular verb, as do compound subjects qualified by \textit{each} or \textit{every}.
\begin{example}
    The long and the short of it is$\ldots$
    
    Bread and butter was all she served.
    
    Give and take is essential to a happy household.
    
    Every window, picture, and mirror was smashed.
\end{example}
A singular subject remains singular even if other nouns are connected to it by \textit{with, as well as, in addition to, except, together with}, and \textit{no less than}.
\begin{example}
    His speech as well as his manner is objectionable.
\end{example}
A linking verb agrees with the number of its subject.
\begin{example}
    What is wanted is a few more pairs of hands.
    
    The trouble with truth is its many varieties.
\end{example}
Some nouns that appear to be plural are usually construed as singular and given a singular verb.
\begin{example}
    Politics is an art, not a science.
    
    The Republican Headquarters is on this side of the tracks.
\end{example}
But
\begin{example}
    The general's quarters are across the river.
\end{example}
In these cases the writer must simply learn the idioms.

The contents of a book is singular.

The contents of a jar may be either singular or plural, depending on what's in the jar - jam or marbles.

\section{Use the proper case of pronoun}
The personal pronouns, as well as the pronoun \textit{who}, change form as they function as subject or object.
\begin{example}
    Will Jane or he be hired, do you think?
    
    The culprit, it turned out, was he.
    
    We heavy eaters would rather walk than ride.
    
    Who knocks?
    
    Give this work to whoever looks idle.
\end{example}
In the last example, \textit{whoever} is the subject of \textit{look idle}; the object of the preposition \textit{to} is the entire clause \textit{whoever looks idle}.

When \textit{who} introduces a subordinate clause, its case depends on its function in that clause.
\begin{example}
    Virgil Soames is the candidate whom we think will win.
    
    $\to$ Virgil Soames is the candidate who we think will win. [We think \emph{he} will win.]
    
    Virgil Soames is the candidate who we hope to elect.
    
    $\to$ Virgil Soames is the candidate whom we hope to elect. [We hope to elect \emph{him}.]
\end{example}
A pronoun in a comparison is nominative if it is the subject of a stated or understood verb.
\begin{example}
    Sandy writes better than I. (Than I write.)
\end{example}
In general, avoid ``understood'' verbs by supplying them.
\begin{example}
    I think Horace admires Jessica more than I.
    
    $\to$ I think Horace admires Jessica more than I do.
    
    Polly loves cake more than me.
    
    $\to$ Polly loves cake more than she loves me.
\end{example}
The objective case is correct in the following examples.
\begin{example}
    The ranger offered Shirley and him some advice on campsites.
    
    They came to meet the Baldwins and us.
    
    Let's talk it over between us, then, you and me.
    
    Whom should I ask?
\end{example}

\begin{example}
    A group of us taxpayers protested.
\end{example}
\textit{Us} in the last example is in apposition to taxpayers, the object of the proposition \textit{of}.

The wording, although grammatically defensible, is rarely apt.

``A group of us protested as taxpayers'' is better, if not exactly equivalent.

%
Use the simple personal pronoun as a subject.
\begin{example}
    Blake and myself stayed home.
    
    $\to$ Blake and I stayed home.
    
    Howard and yourself brought the lunch, I thought.
    
    $\to$ Howard and you brought the lunch, I thought.
\end{example}
The possessive case of pronouns is used to show ownership.

It has 2 forms: the adjectival modifier, \textit{your} hat, and the noun form, a hat \textit{of yours}.
\begin{example}
    The dog has buried 1 of your gloves and 1 of mine in the flower bed.
\end{example}
Gerunds usually require the possessive case.
\begin{example}
    Mother objected to our driving on the icy roads.
\end{example}
A present participle as a verbal, on the other hadn, takes the objective case.
\begin{example}
    They heard him singing in the shower.
\end{example}
The difference between a verbal participle and a gerund is not always obvious, but note what is really said in each of the following.
\begin{example}
    Do you mind me asking a question?
    
    Do you mind my asking a question?
\end{example}
In the 1st sentence, the queried objection is to \textit{me}, as opposed to other members of the group, asking a question.

In the 2nd example, the issue is whether a question may be asked at all.

\section{A participial phrase at the beginning of a sentence must refer to the grammatical subject}
\begin{example}
    Walking slowly down the road, he saw a woman accompanied by 2 children.
\end{example}
The word \textit{walking} refers to the subject to the sentence, not to the woman.

To make it refer to the woman, the writer must recast the sentence.
\begin{example}
    He saw a woman, accompanied by 2 children, walking slowly down the road.
\end{example}
Participial phrases preceded by a conjunction or by a preposition, nouns in apposition, adjectives, and adjective phrases come under the same rule if they begin the sentence.
\begin{example}
    On arriving in Chicago, his friends met him at the station.
    
    $\to$ On arriving in Chicago, he was met at the station by his friends.
    
    A soldier of proved valor, they entrusted him with the defense of the city.
    
    $\to$ A soldier of proved valor, he was entrusted with the defense of the city.
    
    Young and inexperienced, the task seemed easy to me.
    
    $\to$ Young and inexperienced, I thought the task easy.
    
    Without a friend to counsel him, the temptation proved irresistible.
    
    $\to$ Without a friend to counsel him, he found the temptation irresistible.
\end{example}
Sentences violating Rule 11 are often ludicrous:
\begin{example}
    Being in a dilapidated condition, I was able to buy the house very cheap.
    
    Wondering irresolutely what to do next, the clock struck 12.
\end{example}

%------------------------------------------------------------------------------%

\chapter{Elementary Principles of Composition}

\section{Choose a suitable design \& hold to it}
A basic structural design underlies every kind of writing.

Writers will in part follow this design, in part deviate from it, according to their skills, their needs, and the unexpected events that accompany the act of composition.

Writing, to be effective, must follow closely the thoughts of the writer, but not necessarily in the order in which those thoughts occur.

This calls for a scheme of procedure.

In some cases, \textit{the best design is no design}, as with a love letter, which is simply an outpouring, or with a casual essay, which is a ramble.

But in most cases, planning must be a deliberate prelude to writing.

The 1st principle of composition, therefore, is to foresee or determine the shape of what is to come and pursue that shape.

%
A sonnet is built on a 14-line frame, each line containing 5 feet.

Hence, sonneteers know exactly where they are headed, although they may not know how to get there.

Most forms of composition are less clearly defined, more flexible, but all have skeletons to which the writer will bring the flesh and the blood.

The more clearly the writer perceives the shape, the better are the chances of success.

\section{Make the paragraph the unit of composition: 1 paragraph to each topic}
The paragraph is a convenient unit; it serves all forms of literary work.

As long as it holds together, a paragraph may be of any length - a single, short sentence or a passage of great duration.

%
If the subject on which you are writing is of slight extent, or if you intend to treat it briefly, there may be no need to divide it into topics.

Thus, a brief description, a brief book review, a brief account of a single incident, a narrative merely outlining an action, the setting forth of a single idea - any 1 of these is best written in a single paragraph.

After the paragraph has been written, examine it to see whether division will improve it.

%
Ordinarily, however, a subject requires division into topics, each of which should be dealt with in a paragraph.

The object of treating each topic in a paragraph by itself is, of course, to aid the reader.

The beginning of each paragraph is a signal that a new step in the development of the subject has been reached.

%
As a rule, single sentences should not be written or printed as paragraphs.

An exception may be made of sentences of transition, indicating the relation between the parts of an exposition or argument.

%
In dialogue, each speech, even if only a single word, is usually a paragraph by itself; i.e., a new paragraph begins with each change of speaker.

The application of this rule when dialogue and narrative are combined is best learned from examples in well-edited works of fiction.

Sometimes a writer, seeking to create an effect of rapid talk or for some other reason, will elect not to set off each speech in a separate paragraph and instead will run speeches together.

The common practice, however, and the one that serves best in most instances, is to give each speech a paragraph of its own.

%
As a rule, begin each paragraph either with a sentence that suggests the topic or with a sentence that helps the transition.

If a paragraph forms part of a larger composition, its relation to what precedes, or its function as a part of the whole, may need to be expressed.

This can sometimes be done by a mere word or phrase (\textit{again, therefore, for the same reason}) in the 1st sentence.

Sometimes, however, it is expedient to get into the topic slowly, by way of a sentence or 2 of introduction or transition.

%
In narration and description, the paragraph sometimes begins with a concise, comprehensive statement serving to hold together the details that follow.
\begin{example}
    The breeze served us admirably.
    
    The campaign opened with a series of reverses.
    
    The next 10 or 12 pages were filled with a curious set of entries.
\end{example}
But when this device, or any device, is too often used, it becomes a mannerism.

More commonly, the opening sentence simply indicates by its subject the direction the paragraph is to take.
\begin{example}
    At length I thought I might return toward the stockade.
    
    He picked up the heavy lamp from the table and began to explore.
    
    Another flight of steps, and they emerged on the roof.
\end{example}
In animated narrative, the paragraphs are likely to be short and without any semblance of a topic sentence, the writer rushing headlong, event following event in rapid succession.

The break between such paragraphs merely serves the purpose of a rhetorical pause, throwing into prominence some detail of the action.

%
In general, remember that paragraphing calls for a good eye as well as a logical mind.

Enormous blocks of print look formidable to readers, who are often reluctant to tackle them.

Therefore, breaking long paragraphs in 2, even if it is not necessary to do so for sense, meaning, or logical development, is often a visual help.

But remember, too, that firing off many short paragraphs in quick succession can be distracting.

Paragraph breaks used only for show read like the writing of commerce or of display advertising.

Moderation and a sense of order should be the main considerations in paragraphing.

\section{Use the active voice}
48/140

\section{Put statements in positive form}

\section{Use definite, specific, concrete language}

\section{Omit needless words}

\section{Avoid a succession of loose sentences}

\section{Express coordinate ideas in similar form}

\section{Keep related words together}

\section{In summaries, keep to 1 tense}

\section{Place the emphatic words of a sentence at the end}

%------------------------------------------------------------------------------%

\chapter{A Few Matters of Form}

%------------------------------------------------------------------------------%

\chapter{Words \& Expressions Commonly Misused}

%------------------------------------------------------------------------------%

\chapter{An Approach to Style (with a List of Reminders)}

\section{Place yourself in the background}

\section{Write in a way that comes naturally}

\section{Work from a suitable design}

\section{Write with nouns \& verbs}

\section{Revise \& rewrite}

\section{Do not overwrite}

\section{Do not overstate}

\section{Avoid the use of qualifiers}

\section{Do not affect a breezy manner}

\section{Use orthodox spelling}

\section{Do not explain too much}

\section{Do not construct awkward adverbs}

\section{Make sure the reader knows who is speaking}

\section{Avoid fancy words}

\section{Do not use dialect unless your ear is good}

\section{Be clear}

\section{Do not inject opinion}

\section{Use figures of speech sparingly}

\section{Do not take shortcuts at the cost of clarity}

\section{Avoid foreign languages}

\section{Prefer the standard to the offbeat}

%------------------------------------------------------------------------------%

\chapter*{Afterword}

%------------------------------------------------------------------------------%

\printbibliography[heading=bibintoc]
\end{document}